\chapter{结构体与方法}
\label{ch:09}

当一组数据共同描述一个实体时，结构体给字段命名，并让方法围绕同一个状态工作。与用多个平行数组相比，类型能更直接地表达“这些值属于同一条记录”。

本章从计数器的方法调用开始，继续建立泛型模型的具名构造入口，再观察嵌套值的克隆、相等与哈希。字段表示、创建过程和后续合法修改需要一起设计，才能形成可维护的业务模型。

\section{字段与成员方法}

先把计数器看作“一个总数加上修改总数的操作”。结构体给 total 指定类型，成员方法 add 将这项操作放到同一个模型中。方法体用 self 指明当前处理的对象，读代码时便不需要猜测字段属于哪个计数器。

下面让 add 返回 Counter，使调用结果还能继续参与下一次成员调用。先逐次算出总数，再阅读链式形式。返回 self 的具体行为仍由方法类型决定；具名构造入口中的初始化 self 会在下一节另外说明。

下面用计数器说明成员调用和链式返回。返回 self 的含义取决于声明位置：普通方法遵守其返回类型和值、引用规则；具名构造入口另有初始化规则。

\par
\medskip
\begin{lstlisting}[style=chali]
struct Counter
{
    int64 total;

    Counter add(int64 x)
    {
        self.total = self.total + x;
        return self;
    }
};

void main()
{
    Counter c;
    c = c.add(1).add(2).add(3);
    println("{}", c.total);
}
\end{lstlisting}
\par
\medskip

\par
\medskip
\begin{lstlisting}[style=out]
6
\end{lstlisting}
\par
\medskip

\subsection*{代码详解}

Counter 把 total 和 add 方法放在一起，方法内的 self 指当前接收者。初始 total 为 0，每次 add 将参数加到 total，再返回一个 Counter 结果，使后面还能继续调用 add；返回类型不是 int64，因此链式调用中传递的是计数器值。

从左向右跟踪，三次加法分别得到 1、3、6，最外层赋值把最终结果保存回 c，然后打印 6。\texttt{c =\allowbreak{} .\allowbreak{}.\allowbreak{}.\allowbreak{}} 明确保留链式计算的结果；阅读其他方法时也要先看它返回新值、返回自身还是只做修改，不能仅凭点号假定行为相同。

\section{具名构造入口}

static T create 形式的具名入口可以初始化 self 的字段并 return self。调用时使用类型限定名；泛型实参在上下文足够时可以推导。这里的 self 是正在初始化的结果，不是任意静态函数都能访问的全局对象。

\par
\medskip
\begin{lstlisting}[style=chali]
struct Cell<T>
{
    T value;
    static Cell<T> create(T value)
    {
        self.value = value;
        return self;
    }
};
void main()
{
    let cell = Cell::create((int64)42);
    println("{}", cell.value);
}
\end{lstlisting}
\par
\medskip

\par
\medskip
\begin{lstlisting}[style=out]
42
\end{lstlisting}
\par
\medskip

\subsection*{代码详解}

Cell<T> 的字段类型由 T 决定。create 是关联在类型上的静态入口，先把实参写入待返回的 self.value，再返回这个 Cell<T>；它不是要求调用方预先创建 cell 才能调用的普通实例方法。

\texttt{Cell:\allowbreak{}:\allowbreak{}create(\allowbreak{}(\allowbreak{}int64)\allowbreak{}42)\allowbreak{}} 以明确的 int64 实参提供类型信息，编译器据此确定 T，let 再取得整个返回值的类型。打印 cell.value 得到 42。这里推导的是静态类型，不是运行到字段赋值时才动态改变容器的元素类型。

\section{值操作与派生}

结构体的赋值、传参和容器存储沿用第 6 章的值语义。Equal、Hash、Clone 等派生能力必须按支持的指令声明；参与哈希的字段变化后不能把旧索引视为仍然有效。示例把键放入集合，再比较复制前后的嵌套值。

\par
\medskip
\begin{lstlisting}[style=chali]
#derive(Equal, Hash, Clone)
struct BookKey
{
    String label;
    int64 id;
}

#derive(Equal, Hash, Clone)
struct BookRecord
{
    BookKey key;
}

void main()
{
    BookRecord original;
    original.key.label = "left";
    original.key.id = 7;
    BookRecord copied = original.clone();
    Map<BookKey, int64> cache;
    cache.put(original.key, 10);
    println("{} {} {}", original == copied,
            original.hash(original) == copied.hash(copied),
            cache.has(original.key) != 0);
    copied.key.label = "changed";
    println("{} {}", original.key.label, copied.key.label);
}
\end{lstlisting}
\par
\medskip

\par
\medskip
\begin{lstlisting}[style=out]
1 1 1
left changed
\end{lstlisting}
\par
\medskip

\subsection*{代码详解}

BookKey 和外层 BookRecord 都请求相等、哈希与克隆派生，因而可以把内层模型作为 Map 的键，也可以复制整个记录。original 的 label 与 id 设好后，clone 产生 copied，首次输出检查字段相等、相等值的哈希一致，以及刚插入的键确实存在。

随后只改 copied.key.label，第二行分别打印 left 与 changed，说明嵌套模型中的拥有字符串也按其值语义复制。不要把哈希相同反过来当成值相等的证明，也不要在键已用于索引时任意改变决定相等与哈希的字段。本例查询的是原始键，没有展示修改已存键的过程。

\section{组织业务模型}

先确定字段不变量，再决定哪些操作暴露为方法。例如年龄需要范围校验，成绩需要上下界，学生所属班级需要引用仍然存在的班级记录。结构体的类型检查不能代替数据库约束与业务校验。相关的接口抽象放在第 13 章，序列化与反射放在第 20 章。

\subsection*{小结}

结构体把相关字段组织起来，方法控制状态变化，创建入口建立初始状态，派生提供支持范围内的重复值操作。这些机制都继续服从字段类型本来的复制、借用和资源规则。

如果一个模型存在互斥状态，并且不同状态需要不同数据，只堆字段可能产生“哪些字段此时有效”的问题。下一章用枚举把这种选择直接写进类型。

\section{习题}

\begin{enumerate}

\item 给 Cell 增加一个读取方法。

\item 把 Cell 的载荷改成 String。

\item 说明返回 self 的两种场景。

\item 列出学生模型的三条不变量。

\end{enumerate}

\subsection*{参考解答与验收要点}

\begin{enumerate}

\item 返回字段值，按返回类型确定复制或借用。

\item 调用 Cell::create("name") 并核对字段。

\item 构造入口初始化结果；普通方法遵守声明的返回类型。

\item 例如编号唯一、年龄有效、班级存在；在写入边界验证。

\end{enumerate}

下一章：\hyperref[ch:10]{第 10 章}。

\section{番外：木兰的行装}

改写自北朝民歌《木兰诗》。

唧唧复唧唧，木兰改类型。不闻敲键声，唯闻编译鸣。问女何所思，问女何所忆。女亦无所思，女亦无所忆：昨夜见军帖，可汗大点兵；十二份清单，份份不同名。

阿爷无大结构，木兰无统一型。愿为市结构，从此不散拼。东市买 \texttt{horse}，西市买 \texttt{saddle}，南市买 \texttt{bridle}，北市买 \texttt{whip}。四个字段收一处，出发先验能不能骑。

旦辞爷娘去，暮宿边界旁。不闻爷娘唤女声，但闻断言报缺项。弟弟拿绳充马鞍，类型检查不相让；可汗催问能出发否，木兰答曰：“字段全，不等于业务上已经齐。”

万里赴调用，关山度若飞。朔气传日志，寒光照错误。将军百战归，项目十年催。归来见天子，天子问何需。木兰不用加接口，愿借格式化工具，把十二张清单改回一个结构体。

爷娘闻女来，出郭相扶将；阿姊闻妹来，当户改文档；小弟闻姊来，磨刀霍霍向重复字段。同行十二年，不知她的副本早已独立；只因阿爷在家改备注，木兰阵前未受影响。

雄兔脚扑朔，雌兔眼迷离；双兔傍地走，先查模型，安能只凭字段名辨我是公是私？

\subsection*{技术对照}

结构体组织相关字段和行为，但不会自动获得业务有效性。默认值、构造方式和方法应共同维持不变量；派生能力只覆盖对应协议。
